\begin{figure}[htbp]
\centering
\begin{tikzpicture}[x=.56cm,y=.56cm,
  vertice/.style={circle,draw=btinta,fill=white,text=btinta,font=\scriptsize\bfseries,
    inner sep=0pt,minimum size=3.5mm},
  paso/.style={->,>=stealth,draw=btinta,line width=.85pt}]
  \foreach \dx/\izq/\der in {0/1/3,9.7/2/4}{
    \begin{scope}[shift={(\dx,0)}]
      \bcuadriculacero{6}
      \node[vertice] (v1) at ({\izq+.5},4.5) {1};
      \node[vertice] (v2) at ({\izq+.5},1.5) {2};
      \node[vertice] (v3) at ({\der+.5},1.5) {3};
      \node[vertice] (v4) at ({\der+.5},4.5) {4};
      \draw[paso] (v1)--(v2);
      \draw[paso] (v2)--(v3);
      \draw[paso] (v3)--(v4);
      \draw[paso] (v4)--(v1);
      \node[font=\small] at (3,-.73)
        {$R_1\to C_{\izq}\to R_4\to C_{\der}\to R_1$};
    \end{scope}
  }
\end{tikzpicture}
\caption{Two summands of a trace of four incidences on the $6\times6$ grid. Vertex $k$ marks the
intersection where $N_{h_k}$ evaluates $h_k$. On the left, the product is
$h_1(1,1)h_2(4,1)h_3(4,3)h_4(1,3)$. On the right, each column increases by one.
The fourth factor returns to the initial row $R_1$. The sum runs over ordered line indices.
Each formula specifies the measures and multiplicities.}\label{fig:rectangulos}
\end{figure}
